% ============================================================================
% Fichier  : partie6/P06_C25_laboratoire.tex
% Titre    : Le Laboratoire et la Pratique Continue
% Auteur   : MINKA MI NGUIDJOÏ Thierry Emmanuel
% Version  : 1.0 -- novembre 2026
% ============================================================================
% Sources : pratiques d'entraînement en sécurité offensive ; cadre légal des
% laboratoires. NDAP SUKULU : plateforme d'entraînement du cours.
% ============================================================================

\chapter{Le Laboratoire et la Pratique Continue}
\label{chap:labo}
\chaptermark{Le laboratoire et la pratique continue}

\epigraph{%
    « On n'apprend pas à tester en lisant~;\\
    on apprend dans un laboratoire que l'on a le droit d'attaquer. »%
}{Principe du cours}

\niveauChapitre{Tous niveaux}{L'ensemble du cours.}

\begin{boxobjectifs}
\begin{itemize}
    \item Comprendre pourquoi l'entraînement se fait en laboratoire fermé.
    \item Monter un laboratoire personnel légal.
    \item Utiliser les plateformes d'entraînement et les épreuves.
    \item Entretenir ses compétences dans la durée et dans la légalité.
    \item Situer la plateforme NDAP SUKULU du cours.
\end{itemize}
\end{boxobjectifs}

\bigskip

% ============================================================================
\section*{Question directrice : où s'entraîner sans jamais enfreindre la loi~?}
% ============================================================================

Les techniques de ce cours ne se pratiquent que sur des systèmes que l'on a
le droit d'attaquer. La tentation de « s'exercer » sur un site réel est une
faute et une infraction. Ce chapitre montre où et comment s'entraîner
légalement, et présente la plateforme du cours.

\medskip
\textbf{Plan du chapitre.}
\begin{enumerate}
    \item Pourquoi un laboratoire fermé (section~\ref{sec:c25-pourquoi}).
    \item Monter son laboratoire (section~\ref{sec:c25-monter}).
    \item Plateformes et épreuves (section~\ref{sec:c25-plateformes}).
    \item Entretenir ses compétences (section~\ref{sec:c25-entretien}).
    \item La plateforme NDAP SUKULU (section~\ref{sec:c25-ndap}).
\end{enumerate}

\separateur

% ============================================================================
\section{Pourquoi un Laboratoire Fermé}
\label{sec:c25-pourquoi}
% ============================================================================

\begin{boxalert}
\textbf{Violation typique.} S'entraîner sur un site réel « juste pour voir ».

\smallskip
\textbf{Règle~:} tout essai sur un système dont on n'a pas l'autorisation
écrite est une infraction (Chapitre~\ref{chap:droit}), quelles que soient
l'intention et l'absence de dommage. L'entraînement se fait en laboratoire
fermé, isolé du réseau réel.
\end{boxalert}

% ============================================================================
\section{Monter son Laboratoire}
\label{sec:c25-monter}
% ============================================================================

\begin{boxprocedure}
\textbf{Monter un laboratoire personnel légal}
\begin{enumerate}
    \item Un hyperviseur sur sa propre machine~; un réseau virtuel isolé.
    \item Des machines cibles prévues pour l'entraînement (volontairement
          vulnérables, fournies à cet effet).
    \item Aucune connexion de ce réseau au réseau réel ni à l'Internet lors
          des essais offensifs.
    \item Des instantanés pour revenir à l'état initial.
    \item Un journal d'entraînement pour suivre ses progrès.
\end{enumerate}
\end{boxprocedure}

% ============================================================================
\section{Plateformes et Épreuves}
\label{sec:c25-plateformes}
% ============================================================================

Des plateformes en ligne proposent des cibles légales à attaquer et des
épreuves de capture de drapeau. Elles offrent un cadre autorisé~: on y a le
droit explicite de tester. Elles complètent le laboratoire personnel et
permettent de se mesurer.

\begin{tableaucours}{L{4.0cm}L{8.2cm}}{%
    \tableauHeader{Cadre} &
    \tableauHeader{Usage}
}
Laboratoire personnel & Apprendre à son rythme, tout casser sans risque \\
Plateforme d'entraînement en ligne & Cibles variées, légales, mises à jour \\
Épreuve de capture de drapeau & Se mesurer, travailler en équipe, sous
pression \\
Plateforme du cours (NDAP SUKULU) & Fiches, labs et exercices alignés sur ce
cours \\
\end{tableaucours}
\vspace{-0.8em}
\begin{center}{\small\itshape Tableau~25.1. Cadres d'entraînement légaux}\end{center}

% ============================================================================
\section{Entretenir ses Compétences}
\label{sec:c25-entretien}
% ============================================================================

La sécurité offensive évolue vite. Les techniques de fond (reconnaissance,
accès, élévation, latéralité, restitution) restent~; les détails changent.
On entretient ses compétences par une pratique régulière, une veille
mesurée, et le partage au sein d'une communauté.

\begin{boxmetier}{Pentesteur, analyste SOC / Red Team, RSSI}
\textbf{Pratiquer, dans la durée et la légalité.} Le professionnel réserve
un temps régulier au laboratoire, documente ses essais, et ne cède jamais à
la facilité d'un essai « réel » non autorisé~: sa réputation et sa liberté
en dépendent.
\end{boxmetier}

% ============================================================================
\section{La Plateforme NDAP SUKULU}
\label{sec:c25-ndap}
% ============================================================================

Le cours s'appuie sur la plateforme \sigle{NDAP SUKULU}~: pour chaque
chapitre, une fiche d'auto-évaluation~; pour les chapitres pratiques, un lab
qui fait manipuler des pièces (relevés, historiques, journaux) dans un cadre
entièrement fermé. Les dix labs du cours (HE2-LAB-01 à~10) couvrent le
périmètre, l'OSINT, le balayage, la priorisation, l'annuaire, l'élévation,
l'audit de mots de passe, l'historique web, la couverture de détection et le
calcul de gravité.

\separateur

\begin{boxresume}
\textbf{Ce qu'il faut retenir de ce chapitre~:}
\begin{itemize}
    \item \textbf{On s'entraîne} seulement sur ce qu'on a le droit
          d'attaquer.
    \item \textbf{Le laboratoire fermé} est isolé du réseau réel.
    \item \textbf{Plateformes et épreuves} offrent des cibles légales.
    \item \textbf{On entretient} ses compétences par une pratique régulière.
    \item \textbf{NDAP SUKULU} aligne fiches et labs sur le cours.
\end{itemize}
\end{boxresume}

\begin{boxinfo}
\textbf{Points de contrôle de l'élève.} Je sais~: expliquer pourquoi
l'entraînement est en laboratoire fermé~; monter un laboratoire légal~;
citer les cadres d'entraînement~; organiser l'entretien de mes
compétences~; décrire NDAP SUKULU.
\end{boxinfo}

\bigskip

\begin{boxexo}[title={\ding{45}\quad Exercice~25.1. Légalité \niveaubase}]
Pourquoi ne s'entraîne-t-on jamais sur un site réel sans autorisation~?
\end{boxexo}

\begin{boxexo}[title={\ding{45}\quad Exercice~25.2. Isolation \niveaubase}]
Pourquoi le réseau du laboratoire doit-il être isolé du réseau réel~?
\end{boxexo}

\begin{boxexo}[title={\ding{45}\quad Exercice~25.3. Cadres \niveaubase}]
Citez trois cadres d'entraînement légaux.
\end{boxexo}

\begin{boxexo}[title={\ding{45}\quad Exercice~25.4. Instantanés \niveaumoyen}]
À quoi servent les instantanés dans un laboratoire~?
\end{boxexo}

\begin{boxexo}[title={\ding{45}\quad Exercice~25.5. Capture de drapeau \niveaumoyen}]
Qu'apporte une épreuve de capture de drapeau par rapport au laboratoire
personnel~?
\end{boxexo}

\begin{boxexo}[title={\ding{45}\quad Exercice~25.6. Veille \niveaumoyen}]
Comment organiser une veille mesurée sans s'y noyer~?
\end{boxexo}

\begin{boxexo}[title={\ding{45}\quad Exercice~25.7. Journal \niveaumoyen}]
Que consigne-t-on dans un journal d'entraînement~?
\end{boxexo}

\begin{boxexo}[title={\ding{45}\quad Exercice~25.8. NDAP SUKULU \niveaumoyen}]
Décrivez ce que la plateforme du cours apporte à chaque chapitre.
\end{boxexo}

\begin{boxexo}[title={\ding{45}\quad Exercice~25.9. Montage \niveauexpert}]
Décrivez le montage complet d'un laboratoire personnel pour pratiquer les
Parties~III et~IV.
\end{boxexo}

\begin{boxexo}[title={\ding{45}\quad Exercice~25.10. Plan annuel \niveauexpert}]
Proposez votre plan d'entraînement sur un an pour entretenir les compétences
du cours.
\end{boxexo}

\bigskip

\begin{boxinfo}
\textbf{Pour aller plus loin.}
\begin{itemize}
    \item Les dix labs de NDAP SUKULU pour ce cours.
\end{itemize}
\end{boxinfo}

\begin{boxinfo}
\textbf{Transition.} Reste à savoir mesurer la valeur de tout cela~: le
Chapitre~\ref{chap:mesurer} clôt le cours sur la mesure de la sécurité
offensive.
\end{boxinfo}
